\subsection{Engineering Implications and Controlled Replay}

The explanation layer provides a practical bridge between predictive performance and engineering inspection. Attribution-guided masking shows that the highest-ranked TreeSHAP features contain information that is materially relevant to the Stage 2 classifier, while the hierarchical aggregation preserves their provenance from statistical descriptors to physical telemetry channels and uORB topics. This distinction is important because conventional feature importance alone does not directly identify engineering entities that can be inspected. By retaining the telemetry and communication context of each attribution, HS-AeroTS converts model-level explanations into a form that can be interpreted within the PX4 software architecture.

At the same time, prediction relevance should not be equated with physical causality. The weaker and fault-dependent semantic agreement observed at the channel level indicates that features important to the classifier are not necessarily direct indicators of the physical source of a fault. The role of the uORB mapping is therefore not to turn SHAP values into causal fault labels, but to preserve a traceable path from model evidence to architecture-aware inspection candidates. The commit-specific dependency graph is particularly useful in this respect because it associates telemetry evidence with the software revision from which the corresponding topic relationships are derived. The resulting ranking provides an ordered source-tree inspection set, while leaving causal confirmation to subsequent engineering analysis.

Controlled replay provides a stronger test of this interpretation because the mutated source module is known. In this setting, the central question is no longer only whether the model assigns high importance to a telemetry variable, but whether telemetry deviations can retain useful module candidates when compared against a controlled source-level reference. The replay experiments show that matched healthy references can provide such evidence, but also reveal that specificity is at least as important as detection recall. Common-support processing substantially reduced alarms on paired healthy runs while retaining most fault detections, indicating that part of the apparent anomaly evidence in replay originates from differences in observation support rather than from the injected source mutation.

The mutation-specific failures further clarify the current engineering boundary. Persistent EKF2 alarms on healthy replays indicate that normal estimator variation and replay mismatch can dominate residual evidence even in the absence of an injected fault. INAV exhibits a different failure mode: weaker residual responses can remain below a threshold chosen to suppress these normal fluctuations, leading to missed or delayed detections. These two cases expose a fundamental trade-off between sensitivity and specificity in matched-reference replay. They also show why high Top-$K$ module coverage is insufficient on its own: a useful module ranking is only actionable when the upstream detector opens the gate at the appropriate time and does not generate excessive healthy-run alarms.

Taken together, the real-flight explanations and controlled replay results support a diagnostic-triage interpretation of HS-AeroTS. The framework can organize abnormal telemetry into fault-domain evidence and then relate that evidence to a small, architecture-aware set of PX4 modules for inspection. Its current contribution is therefore best understood as reducing the gap between anomaly detection and software inspection, rather than automating software root-cause identification. Achieving the latter would require stronger runtime provenance, broader mutation coverage, and validation on independent source-level faults.
